\documentclass[10pt,letterpaper]{article}
\usepackage[T1]{fontenc}
\usepackage[utf8]{inputenc}
\usepackage[spanish,es-tabla]{babel}
\usepackage{csquotes}
\usepackage{mathpazo}
\usepackage[scaled=0.9]{helvet}
\usepackage{courier}
\usepackage{calc}
\usepackage[top=2cm, bottom=2cm, left=2cm, right=2cm]{geometry}
\usepackage{xcolor}
\usepackage{booktabs}
\usepackage{longtable}
\usepackage{array}
\usepackage{ragged2e}
\usepackage{xurl}
\usepackage[strings]{underscore}
\usepackage{fancyhdr}
\usepackage{sectsty}
\usepackage{tocloft}
\usepackage[colorlinks=true,linkcolor=blue,urlcolor=blue,citecolor=blue,breaklinks=true]{hyperref}
\usepackage{enumitem}

\definecolor{uncpblue}{HTML}{1A237E}
\allsectionsfont{\color{uncpblue}}

\renewcommand{\cftsecleader}{\cftdotfill{\cftsecdotsep}}
\cftpagenumbersoff{section}
\cftpagenumbersoff{subsection}
\cftpagenumbersoff{subsubsection}
\renewcommand{\cftsecfont}{\normalfont}
\renewcommand{\cftsubsecfont}{\normalfont}
\renewcommand{\cftsubsubsecfont}{\normalfont}
\setlength{\parindent}{0pt}
\setlength{\parskip}{4pt}
\setlength{\tabcolsep}{3pt}
\renewcommand{\arraystretch}{0.85}

\pagestyle{fancy}
\fancyhf{}
\fancyhead[L]{\small\textcolor{uncpblue}{\textbf{SGSI-UNCP}}}
\fancyhead[R]{\small\textcolor{uncpblue}{D-SGSI-08}}
\fancyfoot[C]{\thepage}
\renewcommand{\headrulewidth}{0.4pt}
\renewcommand{\footrulewidth}{0.4pt}
\setlength{\headheight}{14pt}

\setlist{nosep,leftmargin=*,after=\vspace{2pt}}
\setlength{\emergencystretch}{2em}

\newcommand{\makecover}{
\begin{titlepage}
\centering
\vspace*{2cm}
{\Huge\bfseries\color{uncpblue} SGSI-UNCP\par}
\vspace{0.7cm}
{\LARGE Guía de Implementación de Controles Críticos (ISO
27002:2022)\par}
\vspace{0.5cm}
{\large Documento Base del Sistema de Gestión\\de Seguridad de la Información\par}
\vspace{1.5cm}
{\small Universidad Nacional del Centro del Perú\par}
\vfill
\end{titlepage}
}

\begin{document}

\makecover
\setcounter{tocdepth}{2}
\RaggedRight
\tableofcontents
\newpage

\RaggedRight
\section{Guía de Implementación de Controles Críticos (ISO
27002:2022)}\label{guuxeda-de-implementaciuxf3n-de-controles-cruxedticos-iso-270022022}

\textbf{Organización:} Universidad Nacional del Centro del Perú (UNCP)
\textbf{Aprobado por:} Oficial de Seguridad y Confianza Digital
\textbf{Basado en:} ISO/IEC 27002:2022 \textbf{Alineamiento:} D-SGSI-05
(Declaración de Aplicabilidad), P-SGSI-03 (Gestión de Accesos),
P-SGSI-02 (Gestión de Incidentes)

\begin{center}\rule{0.5\linewidth}{0.5pt}\end{center}

\subsection{Matriz de Controles de
Vanguardia}\label{matriz-de-controles-de-vanguardia}

Esta tabla detalla los controles priorizados para los activos críticos
de la UNCP (ERP ADESA, Campus Virtual, Datos Personales, infraestructura
cloud), integrando los atributos de la norma y mejoras de automatización
con enfoque Zero Trust.

\subsubsection{1.1 Controles Organizacionales
(A.5)}\label{controles-organizacionales-a.5}

{\setlength{\RaggedRightRightskip}{0pt plus 5em}\small\sloppy \begin{longtable}[]{@{}
  >{\hspace{0pt}\RaggedRight\arraybackslash}p{(\linewidth - 8\tabcolsep) * \real{0.2000}}
  >{\hspace{0pt}\RaggedRight\arraybackslash}p{(\linewidth - 8\tabcolsep) * \real{0.2000}}
  >{\hspace{0pt}\RaggedRight\arraybackslash}p{(\linewidth - 8\tabcolsep) * \real{0.2000}}
  >{\hspace{0pt}\RaggedRight\arraybackslash}p{(\linewidth - 8\tabcolsep) * \real{0.2000}}
  >{\hspace{0pt}\RaggedRight\arraybackslash}p{(\linewidth - 8\tabcolsep) * \real{0.2000}}@{}}
\toprule\noalign{}
\begin{minipage}[b]{\linewidth}\raggedright
ID
\end{minipage} & \begin{minipage}[b]{\linewidth}\raggedright
Control
\end{minipage} & \begin{minipage}[b]{\linewidth}\raggedright
Propósito
\end{minipage} & \begin{minipage}[b]{\linewidth}\raggedright
Atributos (Tipo/CIA/Concepto)
\end{minipage} & \begin{minipage}[b]{\linewidth}\raggedright
Guía de Implementación Mejorada (Zero Trust / Automatización)
\end{minipage} \\
\midrule\noalign{}
\endhead
\bottomrule\noalign{}
\endlastfoot
\textbf{5.1} & Políticas para la seguridad de la información &
Establecer directrices de seguridad & Preventivo / CIA / Gobernar &
\textbf{Documental:} Las políticas se centralizan en el repositorio
documental del SGSI con control de versiones. Se comunican mediante
correo institucional y Moodle. Revisión anual por el CGD. \\
\textbf{5.7} & Inteligencia de Amenazas & Proporcionar conciencia del
entorno de amenazas & Preventivo / CIA / Identificar &
\textbf{Automatización:} Integración de feeds de amenazas (OTX, MISP)
directamente en el SIEM/WAF para bloquear IPs maliciosas de forma
proactiva sin intervención humana. Revisión mensual de nuevas tácticas
MITRE ATT\&CK. \\
\textbf{5.10} & Uso aceptable de la información y de los activos &
Regular el uso correcto de los activos & Preventivo / C / Proteger &
\textbf{Política:} Todo usuario debe aceptar la POL-SGSI-02 al crear su
cuenta institucional. Campañas de concientización trimestrales. \\
\textbf{5.15} & Control de acceso & Limitar el acceso a información y
sistemas & Preventivo / C / Proteger & \textbf{Zero Trust:} Acceso
basado en RBAC + MFA obligatorio para personal. Acceso condicional
(ubicación, dispositivo, riesgo de sesión). Verificación continua
durante toda la sesión. \\
\textbf{5.16} & Gestión de identidades & Gestionar las identidades
digitales & Preventivo / C / Proteger & \textbf{Automatización:} IdM
centralizado (Keycloak/Azure AD) sincronizado con RRHH (SIGA) y Admisión
(ERP ADESA). Creación y baja automatizada de cuentas. Ciclo de vida
completo de la identidad. \\
\textbf{5.17} & Autenticación & Verificar la identidad de los usuarios &
Preventivo / C / Proteger & \textbf{MFA:} MFA obligatorio para todo el
personal administrativo y docente. FIDO2/llaves de hardware para
administradores de sistemas. MFA recomendado para estudiantes. \\
\textbf{5.19} & Seguridad en relaciones con proveedores & Proteger la
información accesible por terceros & Preventivo / CIA / Proteger &
\textbf{Contractual:} Cláusulas de confidencialidad, SLA de seguridad,
derecho a auditoría en todos los contratos de proveedores críticos.
Evaluación precontractual y periódica según POL-SGSI-04. \\
\textbf{5.24} & Gestión de incidentes de seguridad & Asegurar la
respuesta consistente a incidentes & Detectivo / CIA / Responder &
\textbf{CSIRT:} Equipo de respuesta a incidentes con playbooks
documentados. SIEM centralizado con alertas automatizadas. Clasificación
de incidentes según impacto (P-SGSI-02). \\
\textbf{5.30} & Continuidad de la gestión de seguridad & Integrar la
seguridad en la continuidad del negocio & Preventivo / CIA / Recuperar &
\textbf{Arquitectura Híbrida:} RTO \textless{} 1 hora y RPO \textless{}
15 minutos para ERP ADESA mediante replicación sincrónica entre el
Datacenter local y Huawei Cloud. Pruebas de failover semestrales
(P-SGSI-05). \\
\end{longtable}\normalsize}

\subsubsection{1.2 Controles de Personas
(A.6)}\label{controles-de-personas-a.6}

{\setlength{\RaggedRightRightskip}{0pt plus 5em}\small\sloppy \begin{longtable}[]{@{}
  >{\hspace{0pt}\RaggedRight\arraybackslash}p{(\linewidth - 8\tabcolsep) * \real{0.2000}}
  >{\hspace{0pt}\RaggedRight\arraybackslash}p{(\linewidth - 8\tabcolsep) * \real{0.2000}}
  >{\hspace{0pt}\RaggedRight\arraybackslash}p{(\linewidth - 8\tabcolsep) * \real{0.2000}}
  >{\hspace{0pt}\RaggedRight\arraybackslash}p{(\linewidth - 8\tabcolsep) * \real{0.2000}}
  >{\hspace{0pt}\RaggedRight\arraybackslash}p{(\linewidth - 8\tabcolsep) * \real{0.2000}}@{}}
\toprule\noalign{}
\begin{minipage}[b]{\linewidth}\raggedright
ID
\end{minipage} & \begin{minipage}[b]{\linewidth}\raggedright
Control
\end{minipage} & \begin{minipage}[b]{\linewidth}\raggedright
Propósito
\end{minipage} & \begin{minipage}[b]{\linewidth}\raggedright
Atributos
\end{minipage} & \begin{minipage}[b]{\linewidth}\raggedright
Guía de Implementación Mejorada
\end{minipage} \\
\midrule\noalign{}
\endhead
\bottomrule\noalign{}
\endlastfoot
\textbf{6.3} & Concienciación, educación y formación & Asegurar que el
personal conoce sus responsabilidades & Preventivo / C / Proteger &
\textbf{Programa:} Curso anual obligatorio de fundamentos de
ciberseguridad en Moodle. Simulacros de phishing semestrales. Campañas
de concientización mensuales. Programa de capacitación técnica
trimestral (P-SGSI-08). \\
\textbf{6.7} & Trabajo a distancia & Proteger la información accesada
remotamente & Preventivo / CIA / Proteger & \textbf{VPN + MFA:} Acceso
remoto obligatorio mediante VPN institucional con MFA. Prohibición de
uso de redes Wi-Fi públicas para sistemas críticos. Cifrado obligatorio
en dispositivos móviles (POL-SGSI-05). \\
\end{longtable}\normalsize}

\subsubsection{1.3 Controles Físicos
(A.7)}\label{controles-fuxedsicos-a.7}

{\setlength{\RaggedRightRightskip}{0pt plus 5em}\small\sloppy \begin{longtable}[]{@{}
  >{\hspace{0pt}\RaggedRight\arraybackslash}p{(\linewidth - 8\tabcolsep) * \real{0.2000}}
  >{\hspace{0pt}\RaggedRight\arraybackslash}p{(\linewidth - 8\tabcolsep) * \real{0.2000}}
  >{\hspace{0pt}\RaggedRight\arraybackslash}p{(\linewidth - 8\tabcolsep) * \real{0.2000}}
  >{\hspace{0pt}\RaggedRight\arraybackslash}p{(\linewidth - 8\tabcolsep) * \real{0.2000}}
  >{\hspace{0pt}\RaggedRight\arraybackslash}p{(\linewidth - 8\tabcolsep) * \real{0.2000}}@{}}
\toprule\noalign{}
\begin{minipage}[b]{\linewidth}\raggedright
ID
\end{minipage} & \begin{minipage}[b]{\linewidth}\raggedright
Control
\end{minipage} & \begin{minipage}[b]{\linewidth}\raggedright
Propósito
\end{minipage} & \begin{minipage}[b]{\linewidth}\raggedright
Atributos
\end{minipage} & \begin{minipage}[b]{\linewidth}\raggedright
Guía de Implementación Mejorada
\end{minipage} \\
\midrule\noalign{}
\endhead
\bottomrule\noalign{}
\endlastfoot
\textbf{7.1} & Perímetros de seguridad física & Prevenir accesos físicos
no autorizados & Preventivo / CIA / Proteger & \textbf{Datacenter:}
Muros de concreto, puerta blindada, control biométrico + tarjeta, CCTV
24/7 con retención de 90 días, sensor de intrusiones. Campus con
perímetro delimitado y control de acceso vehicular/peatonal
(POL-SGSI-07). \\
\textbf{7.4} & Protección contra amenazas externas y ambientales &
Proteger las instalaciones de desastres & Preventivo / CIA / Proteger &
\textbf{Ambiental:} Detección de incendios (VESDA), extinción por gas
limpio, control de temperatura (18-24°C) y humedad (40-60\%), UPS con 30
min de autonomía + generador eléctrico para 24h. \\
\textbf{7.7} & Escritorio y pantalla limpios & Prevenir la exposición no
autorizada de información & Preventivo / C / Proteger & \textbf{Norma:}
Documentos confidenciales guardados bajo llave al retirarse. Bloqueo
automático de pantalla a los 5 minutos. Prohibición de dejar sesiones
abiertas en estaciones desatendidas. Auditorías visuales
trimestrales. \\
\textbf{7.14} & Eliminación segura de activos & Prevenir la recuperación
de información de activos dados de baja & Preventivo / CIA / Proteger &
\textbf{NIST SP 800-88:} Clear (sobrescritura) para reasignación
interna, Purge (degauss/CE) para baja/donación, Destroy (trituración)
para medios dañados o confidenciales. Certificado de destrucción emitido
(P-SGSI-07). \\
\end{longtable}\normalsize}

\subsubsection{1.4 Controles Tecnológicos
(A.8)}\label{controles-tecnoluxf3gicos-a.8}

{\setlength{\RaggedRightRightskip}{0pt plus 5em}\small\sloppy \begin{longtable}[]{@{}
  >{\hspace{0pt}\RaggedRight\arraybackslash}p{(\linewidth - 8\tabcolsep) * \real{0.2000}}
  >{\hspace{0pt}\RaggedRight\arraybackslash}p{(\linewidth - 8\tabcolsep) * \real{0.2000}}
  >{\hspace{0pt}\RaggedRight\arraybackslash}p{(\linewidth - 8\tabcolsep) * \real{0.2000}}
  >{\hspace{0pt}\RaggedRight\arraybackslash}p{(\linewidth - 8\tabcolsep) * \real{0.2000}}
  >{\hspace{0pt}\RaggedRight\arraybackslash}p{(\linewidth - 8\tabcolsep) * \real{0.2000}}@{}}
\toprule\noalign{}
\begin{minipage}[b]{\linewidth}\raggedright
ID
\end{minipage} & \begin{minipage}[b]{\linewidth}\raggedright
Control
\end{minipage} & \begin{minipage}[b]{\linewidth}\raggedright
Propósito
\end{minipage} & \begin{minipage}[b]{\linewidth}\raggedright
Atributos
\end{minipage} & \begin{minipage}[b]{\linewidth}\raggedright
Guía de Implementación Mejorada
\end{minipage} \\
\midrule\noalign{}
\endhead
\bottomrule\noalign{}
\endlastfoot
\textbf{8.2} & Derechos de acceso con privilegios & Restringir el acceso
a funciones críticas & Preventivo / CI / Proteger & \textbf{Zero Trust
JIT:} Acceso Just-In-Time a cuentas privilegiadas. Los privilegios de
DBA en el ERP ADESA se otorgan por tiempo limitado y bajo aprobación
MFA. Sesiones auditadas en el SIEM. Rotación de credenciales
privilegiadas cada 90 días. \\
\textbf{8.8} & Gestión de vulnerabilidades técnicas & Prevenir la
explotación de vulnerabilidades & Preventivo / CIA / Proteger &
\textbf{Automatización:} Escaneo semanal de vulnerabilidades en todos
los sistemas (locales y cloud). Parches críticos (CVSS \textgreater=
9.0) aplicados en \textless{} 48 horas. Escaneo de contenedores e
imágenes en el registro. Reporte mensual al CGD. \\
\textbf{8.13} & Copias de seguridad & Proteger los datos contra pérdida
& Preventivo / CIA / Proteger & \textbf{Regla 3-2-1:} Respaldo diario de
bases de datos críticas en Huawei Cloud CBR (inmutable). Pruebas de
restauración trimestrales. Cifrado en reposo y en tránsito de todos los
respaldos (POL-SGSI-03). \\
\textbf{8.16} & Seguimiento de actividades (Monitoreo) & Detectar
comportamientos anómalos & Detectivo / CIA / Detectar &
\textbf{UBA/SIEM:} Uso de análisis de comportamiento (User Behavior
Analytics) para detectar accesos anómalos. Integración de logs de Huawei
Cloud, ERP ADESA, Moodle, firewalls y MDM en el SIEM. Alertas en tiempo
real al CSIRT. \\
\textbf{8.20} & Seguridad de redes & Proteger la información en las
redes & Preventivo / CIA / Proteger & \textbf{ZTNA (Microsegmentación):}
Aplicación de Zero Trust Network Access para segmentar lógicamente los
recursos. Un alumno en la WiFi de la facultad no tiene visibilidad de
los servidores de tesorería. Firewall interno entre segmentos académico
y administrativo. 802.1X para autenticación de puertos. \\
\textbf{8.24} & Uso de criptografía & Proteger la confidencialidad e
integridad & Preventivo / CIA / Proteger & \textbf{Cifrado
Generalizado:} TLS 1.3 en todas las comunicaciones externas. Cifrado en
reposo para bases de datos (TDE) y backups. Cifrado de discos
(BitLocker/FileVault) en todos los dispositivos móviles institucionales.
Política de gestión de claves con rotación anual. \\
\textbf{8.25} & Ciclo de vida de desarrollo seguro & Integrar seguridad
en el desarrollo & Preventivo / CI / Proteger & \textbf{DevSecOps:}
Escaneo SAST en cada commit del ERP ADESA. Escaneo DAST trimestral en
aplicaciones web. Análisis de dependencias (SBOM) para identificar
librerías con CVEs. Pruebas de penetración anuales en aplicaciones
críticas (POL-SGSI-01). \\
\textbf{8.31} & Separación de entornos & Evitar contaminación de datos &
Preventivo / CI / Proteger & \textbf{IaC:} Los entornos de desarrollo y
pruebas se despliegan como infraestructura como código
(Terraform/Ansible). Prohibición de uso de datos reales de producción en
entornos no productivos. Datos sintéticos o enmascarados para
pruebas. \\
\textbf{8.32} & Gestión de cambios & Controlar los cambios en los
sistemas & Preventivo / CIA / Proteger & \textbf{ITIL:} Cambios
clasificados como estándar, normal o emergencia (P-SGSI-04). Evaluación
de impacto de seguridad obligatoria para cambios normales. Plan de
rollback documentado. CAB quincenal para cambios mayores. \\
\end{longtable}\normalsize}

\begin{center}\rule{0.5\linewidth}{0.5pt}\end{center}

\subsection{Acciones de Mejora (Superando el
Estándar)}\label{acciones-de-mejora-superando-el-estuxe1ndar}

Para que la UNCP alcance el nivel de madurez deseado en el PGTD y se
diferencie como institución líder en seguridad digital, se sugieren las
siguientes metodologías de mejora continua:

\subsubsection{A. Implementación de SOAR (Security Orchestration,
Automation, and
Response)}\label{a.-implementaciuxf3n-de-soar-security-orchestration-automation-and-response}

\begin{itemize}
\tightlist
\item
  \textbf{Concepto:} Superar el SIEM tradicional automatizando la
  respuesta a incidentes de baja complejidad.
\item
  \textbf{Aplicación UNCP:} Si el SIEM detecta un ataque de fuerza bruta
  contra el VPN, el SOAR instruye automáticamente al Firewall perimetral
  para bloquear la IP atacante, notifica al Oficial de Seguridad por
  Teams y abre un ticket de incidente.
\item
  \textbf{Madurez esperada:} 60\% de los incidentes de nivel Bajo
  gestionados de forma automatizada.
\end{itemize}

\subsubsection{B. Adopción de DevSecOps
Integral}\label{b.-adopciuxf3n-de-devsecops-integral}

\begin{itemize}
\tightlist
\item
  \textbf{Concepto:} La seguridad no es una fase del desarrollo, es
  parte integral del flujo CI/CD.
\item
  \textbf{Aplicación UNCP:} Integrar herramientas de análisis de
  dependencias (OWASP Dependency Check, Snyk) para generar el SBOM del
  Campus Virtual (Moodle) y bloquear el despliegue si se detectan
  librerías con CVEs de severidad crítica.
\item
  \textbf{Madurez esperada:} 100\% de los despliegues escaneados antes
  de pasar a producción.
\end{itemize}

\subsubsection{C. Estrategia de Salida de Nube (Multi-Cloud / Hybrid
Resilience)}\label{c.-estrategia-de-salida-de-nube-multi-cloud-hybrid-resilience}

\begin{itemize}
\tightlist
\item
  \textbf{Concepto:} Evitar el vendor lock-in garantizando que los
  servicios críticos puedan migrar entre proveedores cloud o al entorno
  local.
\item
  \textbf{Aplicación UNCP:} Mantener la capacidad de mover las cargas de
  trabajo críticas (ERP ADESA) desde Huawei Cloud hacia un entorno local
  o una segunda nube pública en menos de 4 horas, mediante el uso de
  contenedores (Kubernetes) y almacenamiento de datos portable.
\item
  \textbf{Madurez esperada:} DRP probado anualmente con conmutación
  exitosa.
\end{itemize}

\subsubsection{D. Programa de Bug Bounty
Interno}\label{d.-programa-de-bug-bounty-interno}

\begin{itemize}
\tightlist
\item
  \textbf{Concepto:} Incentivar la identificación de vulnerabilidades
  por parte de la comunidad universitaria.
\item
  \textbf{Aplicación UNCP:} Programa limitado a estudiantes de
  ingeniería de sistemas y personal de la OTI, con reconocimiento y
  pequeños incentivos por reportes de vulnerabilidades válidos en
  sistemas institucionales.
\end{itemize}

\begin{center}\rule{0.5\linewidth}{0.5pt}\end{center}

\subsection{Mapa de Correlación con la
SoA}\label{mapa-de-correlaciuxf3n-con-la-soa}

{\setlength{\RaggedRightRightskip}{0pt plus 5em}\small\sloppy \begin{longtable}[]{@{}
  >{\hspace{0pt}\RaggedRight\arraybackslash}p{(\linewidth - 4\tabcolsep) * \real{0.3333}}
  >{\hspace{0pt}\RaggedRight\arraybackslash}p{(\linewidth - 4\tabcolsep) * \real{0.3333}}
  >{\hspace{0pt}\RaggedRight\arraybackslash}p{(\linewidth - 4\tabcolsep) * \real{0.3333}}@{}}
\toprule\noalign{}
\begin{minipage}[b]{\linewidth}\raggedright
Control de esta Guía
\end{minipage} & \begin{minipage}[b]{\linewidth}\raggedright
ID en la SoA (D-SGSI-05)
\end{minipage} & \begin{minipage}[b]{\linewidth}\raggedright
Estado Objetivo 2027
\end{minipage} \\
\midrule\noalign{}
\endhead
\bottomrule\noalign{}
\endlastfoot
5.1, 5.7, 5.10, 5.15, 5.16, 5.17 & Organizacionales (A.5) &
Implementado \\
5.19, 5.24, 5.30 & Organizacionales (A.5) & Implementado \\
6.3, 6.7 & Personas (A.6) & Implementado \\
7.1, 7.4, 7.7, 7.14 & Físicos (A.7) & Implementado \\
8.2, 8.8, 8.13, 8.16, 8.20 & Tecnológicos (A.8) & Implementado \\
8.24, 8.25, 8.31, 8.32 & Tecnológicos (A.8) & Implementado \\
\end{longtable}\normalsize}

\begin{center}\rule{0.5\linewidth}{0.5pt}\end{center}

\subsection{Restricción de Confidencialidad y
Cumplimiento}\label{restricciuxf3n-de-confidencialidad-y-cumplimiento}

Todas las recomendaciones contenidas en esta guía están diseñadas para
cumplir con la \textbf{Ley N° 29733} (Protección de Datos Personales),
el \textbf{D.L. N° 1412} (Ley de Gobierno Digital), el \textbf{D.S. N°
016-2024-JUS} (Reglamento LPDP) y el \textbf{D.S. N° 029-2021-PCM}
(Reglamento de Gobierno Digital). Este documento es para uso interno de
la UNCP; su distribución externa requiere la anonimización de activos
específicos y la autorización del Oficial de Seguridad.

\begin{center}\rule{0.5\linewidth}{0.5pt}\end{center}

\subsection{Documentos Relacionados}\label{documentos-relacionados}

{\setlength{\RaggedRightRightskip}{0pt plus 5em}\small\sloppy \begin{longtable}[]{@{}ll@{}}
\toprule\noalign{}
Código & Nombre \\
\midrule\noalign{}
\endhead
\bottomrule\noalign{}
\endlastfoot
D-SGSI-05 & Declaración de Aplicabilidad (SoA) \\
P-SGSI-03 & Gestión de Identidades y Control de Acceso \\
P-SGSI-02 & Marco de Respuesta a Incidentes (CSIRT) \\
P-SGSI-05 & Resiliencia y Continuidad en Nube Híbrida \\
POL-SGSI-01 & Política de Desarrollo Seguro \\
POL-SGSI-07 & Política de Seguridad Física \\
R-SGSI-01 & Inventario de Activos de Información \\
\end{longtable}\normalsize}

\begin{center}\rule{0.5\linewidth}{0.5pt}\end{center}

\end{document}
